\subsection*{(a) Impacto do tamanho da janela sem perda}

No GBN, o throughput passou de 181,9\,kbit/s (SAW, $N=1$) para 529,4, 723,6 e 1\,226,2\,kbit/s com $N = 4$, 8 e 16 (2,9$\times$, 4,0$\times$ e 6,7$\times$ o SAW). De 4 para 8 o ganho foi de 37\% e de 8 para 16, de 69\%. No SR, o throughput passou de 181,9\,kbit/s (SAW, $N=1$) para 374,1, 1\,058,4 e 2\,132,2\,kbit/s com $N = 4$, 8 e 16 (2,1$\times$, 5,8$\times$ e 11,7$\times$ o SAW). De 4 para 8 o ganho foi de 183\% e de 8 para 16, de 101\%.

\par Sem perda, GBN e SR executam o mesmo caminho (pedem, recebem em ordem e confirmam, sem retransmissões), então diferenças entre eles nessa coluna vêm das condições da rede em cada execução, e não do protocolo. Com $N$ segmentos em trânsito, chegam até $N$ segmentos por RTT, e o teto ideal é $N \times L/\mathrm{RTT}$ (com RTT mediano de 18,68\,ms, 877, 1\,755 e 3\,509\,kbit/s para $N = 4$, 8 e 16). Esse crescimento linear vale enquanto a janela é menor que o produto banda $\times$ atraso do caminho. Quando $N \times L$ passa a cobrir todo o RTT, o canal fica sempre ocupado e o limite passa a ser o gargalo (a banda do enlace ou a taxa de resposta do servidor e do cliente), não mais a janela.

\par Nos dados, o throughput \textbf{não se estabilizou até $N = 16$}: cada vez que a janela dobrou, a vazão ainda subiu de forma expressiva, sinal de que o produto banda $\times$ atraso do caminho não foi atingido. A estabilização só ocorreria com janelas maiores. A maior vazão medida (2\,132,2\,kbit/s) equivale a $\approx 9,7$ segmentos entregues por RTT mediano, abaixo dos 16 da janela, porque parte do tempo a janela fica parada esperando o \texttt{DATA} mais lento. A estabilização ocorreria na janela em que $N \times L/\mathrm{RTT}$ alcançasse a banda disponível. Ver Figura~\ref{fig:janela}.

\par Dois cuidados na leitura desses números. Primeiro, cada cenário foi executado uma vez, e o RTT mediano variou de 18,7 a 38,0\,ms entre as execuções. Medida contra o ideal $N \times L/\mathrm{RTT}$ da própria execução, a eficiência ficou entre 55\% e 83\%, porque os picos de RTT seguram a janela e o \texttt{DATA} mais lento atrasa o deslizamento. Segundo, mesmo sem perda simulada houve \textbf{perda natural} na rede (retransmissões com \texttt{loss\_simulated}$=0$): no GBN com $N=16$, 1 timeout (16 \texttt{REQ}s reenviados) custou $\approx 1$\,s de 6,8\,s; sem ele, a vazão seria $\approx 1\,436$\,kbit/s; no SR com $N=4$, 3 timeouts (3 \texttt{REQ}s reenviados) custaram $\approx 3$\,s de 22,4\,s; sem eles, a vazão seria $\approx 432$\,kbit/s. Numa transferência de poucos segundos, um único timeout de 1\,s pesa muito na média.

\subsection*{(b) Retransmissões por erro com 10\% de perda e janela 16}

No GBN, cada timeout da base reenvia a janela inteira, somando $+1$ na coluna \texttt{retransmissions} de até 16 segmentos consecutivos. Separando esses blocos no CSV, houve 243 timeouts e 3878 \texttt{REQ}s reenviados: \textbf{16,0 segmentos retransmitidos por erro}, praticamente a janela inteira. No SR, cada timeout reenvia um único \texttt{REQ}: houve 234 retransmissões para 193 segmentos com \texttt{loss\_simulated}$=1$, ou \textbf{1,21 por erro}. O valor fica um pouco acima de 1 porque o próprio reenvio também pode ser descartado. No total, o GBN reenviou \textbf{16,6$\times$} mais \texttt{REQ}s que o SR. O GBN também marcou mais segmentos com perda simulada (469 contra 193), porque recebe muito mais \texttt{DATA}s, e cada reenvio da janela é exposto de novo aos 10\% de descarte (Figura~\ref{fig:segmentos}).

\par Um exemplo do CSV: no GBN, o primeiro timeout ocorreu na base \texttt{seq=1}, e as linhas 1 a 16 (o bloco de 16 segmentos daquela janela) passaram a ter \texttt{retransmissions}$\geq1$, todos reenviados por causa de um único \texttt{DATA} perdido. No SR, o primeiro segmento retransmitido foi o \texttt{seq=1}, e nenhum dos três vizinhos de cada lado foi reenviado.

\subsection*{(c) Por que o GBN sofre mais que o SR com alta perda}

Com 10\% de perda (Tabela~\ref{tab:comparativa}), com $N=4$, o GBN fez 1\,010 retransmissões e levou 306,4\,s (27,4\,kbit/s), contra 250 retransmissões e 231,2\,s (36,3\,kbit/s) do SR; com $N=8$, o GBN fez 1\,912 retransmissões e levou 263,5\,s (31,8\,kbit/s), contra 221 retransmissões e 152,2\,s (55,1\,kbit/s) do SR; com $N=16$, o GBN fez 3\,878 retransmissões e levou 251,5\,s (33,4\,kbit/s), contra 234 retransmissões e 109,0\,s (77,0\,kbit/s) do SR.

Com 10\% de perda, aumentar a janela de 4 para 16 mudou o throughput do GBN em +22\% e o do SR em +112\%.

\par A diferença vem do receptor. O GBN não guarda segmentos fora de ordem: quando o \texttt{DATA} da base se perde, os $N-1$ \texttt{DATA}s seguintes, que chegaram corretamente, são descartados, e todos são pedidos de novo após o timeout. Com perda $p$, a chance de uma janela de $N$ segmentos passar sem nenhum erro é $(1-p)^N$, que vale 19\% para $p=10\%$ e $N=16$. Quase toda janela tem um erro, e cada erro custa até $N$ reenvios. Quanto maior a janela, mais trabalho já feito é jogado fora a cada erro. O SR guarda os segmentos fora de ordem no buffer e reenvia só o que faltou: o custo de um erro é um \texttt{REQ}, independentemente de $N$. Enquanto espera o timeout do segmento perdido, ele continua recebendo o resto da janela, e por isso janelas maiores continuam ajudando o SR mesmo com perda.
